\section{Hardness of Graph Partitioning with Demands}\label{sec:gpd-hardness}

We next provide an approximation-hardness result for \textsc{Graph Partitioning with Demands} based on the Small-Set Expansion Hypothesis. The \emph{Small-Set Expansion Hypothesis} states that, for every constant $\zeta>0$, there is a constant $\delta>0$ for which it is NP-hard, given a regular graph, to distinguish between the existence of a vertex set of normalized size $\delta$ and expansion at most $\zeta$, and the condition that every vertex set of normalized size $\delta$ has expansion at least $1-\zeta$.

For the weighted graph $X$ in the reduction, define the normalized vertex measure
\[
    \mu(S)=\frac{\operatorname{vol}_X(S)}{\operatorname{vol}_X(U)}.
\]
Thus $\mu(U)=1$ and $\mu(u)$ is the stationary probability of stepping into $u$ under the random walk on $X$. Since multiplying all edge weights by the same constant changes neither expansion nor approximation ratios, we normalize $\operatorname{vol}_X(U)=1$. With this normalization, the expansion of a nonempty set $S$ is
\[
    \Phi_X(S)=\frac{c\br{S,U\setminus S}}{\mu(S)}.
\]

We also recall the Gaussian expansion profile used in the source theorem. For $r\in[0,1)$ and $\alpha\in(0,1)$, let
\[
    \begin{pmatrix}Z\\Z'\end{pmatrix}
    \sim\mathcal{N}\left(
        \begin{pmatrix}0\\0\end{pmatrix},
        \begin{pmatrix}1&r\\r&1\end{pmatrix}
    \right).
\]
Thus $Z$ and $Z'$ are standard Gaussian random variables with correlation $r$. Let $t_\alpha\in\mathbb{R}$ satisfy $\Pr[Z\geq t_\alpha]=\alpha$. We define the expansion profile of the Gaussian noise graph $\mathcal{G}(r)$ by
\[
    \Phi_{\mathcal{G}(r)}(\alpha)
    =\Pr\left[Z'<t_\alpha\mid Z\geq t_\alpha\right]
    =1-\frac{\Pr[Z\geq t_\alpha,\,Z'\geq t_\alpha]}{\alpha}.
\]

We first invoke the following result of Raghavendra, Steurer, and Tulsiani~\cite[Theorem~3.5]{ReductionsBetweenExpansionProblems}.

\begin{theorem}[Raghavendra-Steurer-Tulsiani]\label{thm:rst-global-expansion}
    For every integer $q\geq2$ and every $\varepsilon,\gamma>0$, assuming the Small-Set Expansion Hypothesis, it is hard to distinguish between the following cases for a weighted graph $X=(U,F,c)$:
    \begin{description}
        \item[YES:] There is a partition $\cP=\{S_1,\ldots,S_q\}$ of $U$ such that, for every $i\in[q]$, 
            $\mu(S_i)=\frac{1}{q}$ and $\Phi_X(S_i)\leq\varepsilon+o(\varepsilon)$.
        \item[NO:] Every $S\subseteq U$ satisfies $
            \Phi_X(S)
            \geq
            \Phi_{\mathcal{G}(1-\varepsilon/2)}(\mu(S))
            -\frac{\gamma}{\mu(S)}$.
    \end{description}
\end{theorem}

We derive from Theorem~\ref{thm:rst-global-expansion} the form of the gap used in our reduction.

\begin{lemma}\label{lem:globalized-long-code-gap}
    Fix an integer $q\geq2$ and a constant $b\in(0,1/2]$. Assuming the Small-Set Expansion Hypothesis, for every sufficiently small $\varepsilon>0$ it is hard to distinguish between the following cases for a weighted graph $X=(U,F,c)$:
    \begin{description}
        \item[YES:] There is a partition $\cP=\brc{S_1,\ldots,S_q}$ of $U$
        such that, for every $i\in[q]$, $
            \mu(S_i)=\frac{1}{q}$ and $
            \Phi_X(S_i)\leq K\cdot\varepsilon,$
        where $K>0$ is a universal constant.
        \item[NO:] For every $S\subseteq U$, where $b\leq\mu(S)\leq\frac{1}{2}$, $\Phi_X(S)\geq a_b\cdot \sqrt{\varepsilon},$
        where $a_b>0$ depends only on $b$.
    \end{description}
\end{lemma}
\begin{proof}
    In the YES case of Theorem~\ref{thm:rst-global-expansion}, the sets $S_1,\ldots,S_q$ have total measure $1$ and therefore form a partition of $U$. Moreover, $\varepsilon+o(\varepsilon)\leq K\cdot\varepsilon$ for a universal constant $K$ and all sufficiently small $\varepsilon$.

    Consider now the NO case. If $\alpha\in[b,1/2]$, then $t_\alpha$ ranges over a bounded interval depending only on $b$. The standard estimate for the probability that two highly correlated Gaussian variables lie on opposite sides of such a threshold gives $
        \Phi_{\mathcal{G}(1-\varepsilon/2)}(\alpha)
        \geq a'_b\cdot\sqrt{\varepsilon}$
    for some constant $a'_b>0$ and all sufficiently small $\varepsilon$. Indeed, the standard deviation of $Z'-Z$ is $\sqrt{\varepsilon}$, while the Gaussian density at $t_\alpha$ is bounded above and below by positive constants depending only on $b$. Taking $\gamma\leq b\cdot a'_b\cdot\sqrt{\varepsilon}/2$ yields
    \[
        \Phi_X(S)\geq\frac{a'_b}{2}\cdot\sqrt{\varepsilon}
    \]
    whenever $b\leq\mu(S)\leq1/2$. The claim follows with $a_b=a'_b/2$.
\end{proof}

\subsection{The reduction}

For simplicity, we carry out the reduction with rational weights. In order for the capacities and demands to be natural numbers, we may always multiply all of them by a common denominator. This does not affect feasibility or any approximation ratio, and hence does not affect the analysis.

A demand graph is \emph{bipartite multiplicative} if its vertex set has a bipartition $L\mathbin{\dot\cup}R$, its demand edges are all pairs with one endpoint in each side, and there is a potential $\pi$ such that $w(uv)=\pi(u)\cdot \pi(v)$ for every demand edge $uv$. Our intention is to show a reduction which establishes the hardness of graph partitioning with bipartite multiplicative demands.
This class is closed under taking induced subgraphs.


    For every $u\in U$, introduce two vertices $u^{\mathrm L}$ and $u^{\mathrm R}$, assign $
        \pi\br{u^{\mathrm L}}=\pi\br{u^{\mathrm R}}=\mu(u),$
    and let $
        \widetilde U
        =\brc{u^{\mathrm L},u^{\mathrm R}\colon u\in U}.$
    On $\widetilde U$, define the bipartite multiplicative demand graph to have edge set $
        D_\times
        =\brc{u^{\mathrm L}v^{\mathrm R}\colon u,v\in U\times U}$ and demands $
        w\br{u^{\mathrm L},v^{\mathrm R}}
        =\pi\br{u^{\mathrm L}}\cdot \pi\br{v^{\mathrm R}}.$
    For $S\subseteq U$, write $
        \widetilde S=\brc{u^{\mathrm L},u^{\mathrm R}\colon u\in S}.$
\begin{lemma}\label{lem:multiplicative-demands-encode-measure}
    For any $S\subseteq U$, we have, $
        {w\br{\widetilde S}}
        =\mu(S)^2$. In particular, $w\br{\widetilde U}=1$.
\end{lemma}
\begin{proof}
    Every pair $(u,v)\in S\times S$ indexes the unique demand edge $u^{\mathrm L}v^{\mathrm R}$. Therefore,
    \[
        w\br{\widetilde S}
        =\sum_{(u,v)\in S\times S}\pi\br{u^{\mathrm L}}\cdot \pi\br{v^{\mathrm R}}
        =\left(\sum_{u\in S}\mu(u)\right)^2
        =\mu(S)^2.
    \]
    The same calculation for $S=U$ gives $w\br{\widetilde U}=1$.
\end{proof}

We now turn $X$ into a supply graph $\widetilde X$ on $\widetilde U$. For every $uv\in F$, add the supply edge $u^{\mathrm L}v^{\mathrm L}$ with capacity $c(uv)$. For every $u\in U$, also add the \emph{consistency edge} $u^{\mathrm L}u^{\mathrm R}$ with capacity $1$. Thus, unless a consistency edge is cut, the two copies of every vertex belong to the same component. Together with the previously defined bipartite multiplicative demand graph on $\widetilde U$, this completes the construction of the instance $(\widetilde X,D_\times)$.

\subsection{Completeness}
\begin{lemma}\label{lem:sse-gpd-completeness}
    Fix $\rho\in(0,1)$ and set $
        q=\max\left\{2,\left\lceil\rho^{-1/2}\right\rceil\right\}. $
    In the YES case,
    \[
        \OPTdem{\widetilde X}{\rho}\leq \frac{K}{2}\cdot\varepsilon.
    \]
\end{lemma}
\begin{proof}
    For every $i\in[q]$, define $
        \widetilde S_i
        =\brc{u^{\mathrm L},u^{\mathrm R}\colon u\in S_i}.$
    Since $S_1,\ldots,S_q$ partition $U$, the sets $\widetilde S_1,\ldots,\widetilde S_q$ partition $\widetilde U$. Let $C_{\mathrm{yes}}$ consist of the supply edges $u^{\mathrm L}v^{\mathrm L}$ for which $uv\in F$ has endpoints in distinct parts of the original partition. Thus all consistency edges $u^{\mathrm L}u^{\mathrm R}$ are retained. Every crossing edge of $X$ is counted twice in the sum of the boundaries, and hence
    \[
        c\br{C_{\mathrm{yes}}}
        =\frac{1}{2}\sum_{i=1}^{q}c\br{S_i,U\setminus S_i}.
    \]
    By Lemma~\ref{lem:globalized-long-code-gap},
    \[
        c\br{C_{\mathrm{yes}}}
        \leq\frac{1}{2}\sum_{i=1}^{q}\mu(S_i)\cdot K\cdot\varepsilon
        =\frac{K\cdot\varepsilon}{2}.
    \]

    Every component $\widetilde H$ of $\widetilde X\setminus C_{\mathrm{yes}}$ is a union of left-right pairs corresponding to a set $H\subseteq S_i$ for some $i$. Therefore, by Lemma~\ref{lem:multiplicative-demands-encode-measure},
    \[
        \frac{w\br{\widetilde H}}{w\br{\widetilde U}}
        =\mu(H)^2
        \leq\mu(S_i)^2
        =\frac{1}{q^2}
        \leq\rho.
    \]
    Thus $C_{\mathrm{yes}}$ is $\rho$-feasible and has cost at most $\frac{K\cdot\varepsilon}{2}$.
\end{proof}

\subsection{Soundness}
We start by proving a numerical lemma that will be useful in the soundness analysis.
\begin{lemma}\label{lem:group-small-components}
    Let $p_1,\ldots,p_m\geq0$ satisfy $
        \sum_{i=1}^{m}p_i=1$ and $
        p_i\leq\sqrt{\eta}<1.$
    Define $b_\eta=\min\left\{\frac{1}{4},1-\sqrt{\eta}\right\}$.
    Then, there is a set $I\subseteq[m]$ such that
    \[
        b_\eta\leq\sum_{i\in I}p_i\leq\frac{1}{2}.
    \]
\end{lemma}
\begin{proof}
    If some $p_i>1/2$, take all indices except $i$. Their total is at least $1-\sqrt{\eta}\geq b_\eta$ and is less than $1/2$.
    Otherwise, we have $p_i\leq1/2$. If some $p_i\geq b_\eta$, take that index. If all $p_i<b_\eta$, add indices until their sum first reaches $b_\eta$. The resulting sum is less than $2b_\eta\leq1/2$.
\end{proof}

\begin{lemma}\label{lem:sse-gpd-soundness}
    In the NO case, every $\eta$-feasible supply-edge set $C$ satisfies $
        c\br{C}\geq\beta_\eta\cdot\sqrt{\varepsilon}
    $
    for some constant $\beta_\eta>0$ depending only on $\eta$, provided that $\varepsilon$ is sufficiently small.
\end{lemma}
\begin{proof}
    Set $\beta_\eta=b_\eta\cdot a_{b_\eta}$. We may assume that $\beta_\eta\cdot\sqrt{\varepsilon}\leq1$. If $C$ contains a consistency edge, then $c\br{C}\geq1\geq\beta_\eta\cdot\sqrt{\varepsilon}$, as required.

    Suppose now that no consistency edge is cut. Let $\widetilde H_1,\ldots,\widetilde H_m$ be the components of $\widetilde X\setminus C$. Every $\widetilde H_i$ is a union of left-right pairs. Let $H_i\subseteq U$ be their projection and put $p_i=\mu(H_i)$. The sets $H_1,\ldots,H_m$ partition $U$. By $\eta$-feasibility and Lemma~\ref{lem:multiplicative-demands-encode-measure},
    \[
        p_i^2 = \mu(H_i)^2
        =\frac{w\br{\widetilde H_i}}{w\br{\widetilde U}}
        \leq\eta,
    \]
    and hence $p_i\leq\sqrt{\eta}$. Apply Lemma~\ref{lem:group-small-components} and let $
        S=\bigcup_{i\in I}H_i.$
    Then $b_\eta\leq\mu(S)\leq1/2$. Every original-copy supply edge corresponding to an edge of $X$ crossing $\br{S,U\setminus S}$ joins two distinct components of $\widetilde X\setminus C$ and therefore belongs to $C$. Consequently, $
        c\br{C}\geq c\br{S,U\setminus S},
    $ so by Lemma~\ref{lem:globalized-long-code-gap},
    \[
        c\br{S,U\setminus S}
        =\mu(S)\cdot\Phi_X(S)
        \geq b_\eta\cdot a_{b_\eta}\sqrt{\varepsilon}
        =\beta_\eta\cdot\sqrt{\varepsilon}.
    \]
    which completes the proof.
\end{proof}
\subsection{Hardness of approximation}
\begin{theorem}\label{thm:sse-gpd-bicriteria-hardness}
    Fix constants $0<\rho<\eta<1$. Then, for any sufficiently small $\varepsilon>0$, a universal constant $L$ and a constant $\beta_\eta>0$ dependent only on $\eta$ such that it is SSEH-hard to distinguish between the following two cases:
    \begin{itemize}
        \item \textbf{YES case:} $\OPTdem{\widetilde X}{\rho}\leq L\cdot\varepsilon$.
        \item \textbf{NO case:} Every $\eta$-feasible supply-edge set $C$ satisfies $c\br{C}\geq\beta_\eta\cdot\sqrt{\varepsilon}$.
    \end{itemize}
    Consequently, assuming the Small-Set Expansion Hypothesis, for any constant $A>0$ there is no polynomial-time $(A,\eta)$-bicriteria approximation for \textsc{Graph Partitioning with Demands}, even for bipartite multiplicative demands.
\end{theorem}
\begin{proof}
    Let $L=\frac{K}{2}$. 
    In the YES case, Lemma~\ref{lem:sse-gpd-completeness} gives $\OPTdem{\widetilde X}{\rho}\leq L\cdot\varepsilon$. 
    In the NO case, Lemma~\ref{lem:sse-gpd-soundness} says that every $\eta$-feasible cut has cost at least $\beta_\eta\cdot\sqrt{\varepsilon}$. 
    Now, suppose that such an algorithm exists. Choose $\varepsilon>0$ sufficiently small that all preceding statements apply and
    \[
        A\cdot L\cdot\varepsilon<\beta_\eta\cdot\sqrt{\varepsilon}.
    \]
    Equivalently, it is enough to take
    \[
        \varepsilon<\left(\frac{\beta_\eta}{A\cdot L}\right)^2.
    \]
    Thus the algorithm distinguishes the YES and NO instances of Lemma~\ref{lem:globalized-long-code-gap}, contradicting the Small-Set Expansion Hypothesis.
\end{proof}

The hardness also transfers to the \textsc{Generalized Conductance} problem, as formalized in the following corollary.

\begin{corollary}\label{cor:sse-generalized-conductance-hardness}
    It is SSEH-hard to approximate \textsc{Generalized Conductance} within any constant factor, even for bipartite multiplicative demands.
\end{corollary}
\begin{proof}
    Theorem~\ref{thm:gpd-extension} certifies that a constant-factor approximation for \textsc{Generalized Conductance} would yield a constant-factor bicriteria approximation for \textsc{Graph Partitioning with Demands}, contradicting Theorem~\ref{thm:sse-gpd-bicriteria-hardness}.
\end{proof}
